\documentclass[11pt]{amsart}

\usepackage[margin=1.15in]{geometry}
\usepackage{amsmath,amssymb,amsthm,mathtools}
\usepackage{booktabs}
\usepackage[font=small,width=0.85\textwidth]{caption}
\usepackage{enumitem}
\usepackage[hidelinks]{hyperref}

\newtheorem{theorem}{Theorem}[section]
\newtheorem{proposition}[theorem]{Proposition}
\newtheorem{lemma}[theorem]{Lemma}
\newtheorem{corollary}[theorem]{Corollary}
\newtheorem{conjecture}[theorem]{Conjecture}
\theoremstyle{definition}
\newtheorem{definition}[theorem]{Definition}
\newtheorem{observation}[theorem]{Observation}
\theoremstyle{remark}
\newtheorem{remark}[theorem]{Remark}

\newcommand{\F}{\mathbb{F}}
\newcommand{\Q}{\mathbb{Q}}
\newcommand{\R}{\mathbb{R}}
\newcommand{\C}{\mathcal{C}}
\newcommand{\conj}[1]{\overline{#1}}
\newcommand{\assoc}[3]{[#1,#2,#3]}
\DeclareMathOperator{\rank}{rank}
\DeclareMathOperator{\tr}{tr}
\DeclareMathOperator{\Ort}{O}
\DeclareMathOperator{\SO}{SO}
\DeclareMathOperator{\spn}{span}

\title[Square-root-free associator defects]{Square-root-free associator defects in octonion algebras over finite fields, with a mutation-tested pentagon oracle}

\author{[Author names withheld for review]}

\date{October 1, 2026}

\begin{document}

\begin{abstract}
For an octonion algebra $\C$ over a field of characteristic different from~$2$
and elements $x,y$ of nonzero norm, we study the defect operator
$\Delta(x,y)=L_xL_yL_{xy}^{-1}$, a square-root-free replacement for the
normalized operator $U_xU_yU_{xy}^{-1}$, $U_x=L_x/\sqrt{N(x)}$, which is
unavailable over finite fields. We prove that $\Delta(x,y)$ is a
special orthogonal transformation of $(\C,N)$ that fixes pointwise the unital
subalgebra generated by $x$ and $y$. When that subalgebra is a nondegenerate
quaternion subalgebra $Q$, the defect is therefore an element of
$\SO(Q^\perp)$, a four-dimensional rotation group, rather than merely of the
stabilizer of~$xy$. The determinant statement rests on the polynomial identity
$\det L_x=N(x)^4$. A seeded census over the split octonions modulo
$\ell\in\{3,5,7,11\}$ supports a conjectural rank law
$\rank(\Delta-I)\in\{0,2,4\}$ determined by $\dim Q$. We prove its first case. We then describe an
executable Lean~4 oracle for the five-term associator (Teichm\"uller) identity,
organized as a pentagon of parenthesizations. Mutation testing shows that the
closure of that identity is a weak test of an implementation: it is invariant
under a global sign change of the associator and under passage to the opposite
algebra, and only replay against an independent implementation detects the
latter. All algebraic results are elementary consequences of classical theory;
the contribution is the exact finite formulation, its structural consequences,
and a reproducible certification methodology.
\end{abstract}

\subjclass[2020]{Primary 17D05; Secondary 17A75, 11E57, 68V15}
\keywords{octonions, composition algebras, finite fields, associator, Teichm\"uller identity, mutation testing, Lean}

\maketitle

\section{Introduction}\label{sec:intro}

Left multiplication by a nonzero element $x$ of a real division octonion
algebra is a similitude of the norm form. After normalization,
$U_x=L_x/\sqrt{N(x)}$ lies in $\SO(8)$. Because the octonions are not
associative, $x\mapsto U_x$ is not a homomorphism, and the failure is measured
by the operator
\begin{equation}\label{eq:real-defect}
  D(x,y)=U_xU_yU_{xy}^{-1}\in\SO(8),
\end{equation}
which is the identity exactly when $x(yz)=(xy)z$ for every $z$. Operators of
this kind are natural observables for word problems in the matrix group
generated by the $U_x$. Over a finite field, however, the normalization
$\sqrt{N(x)}$ need not exist, and~\eqref{eq:real-defect} is not even defined.

This paper records what survives without square roots and what can be
certified about it by exact computation. Our results are as follows.

\begin{enumerate}[label=(\roman*)]
\item The square-root-free operator $\Delta(x,y)=L_xL_yL_{xy}^{-1}$ coincides
  with $D(x,y)$ over $\R$ and is defined over every field of characteristic
  $\neq2$. It is an isometry fixing $xy$, and it is trivial exactly when the
  associator $[x,y,\cdot\,]$ vanishes (Proposition~\ref{prop:basic}).
\item $\Delta(x,y)$ fixes pointwise the unital subalgebra $A$ generated by $x$
  and $y$, preserves $W^\perp$ for every subspace $W\subseteq A$, and has
  determinant~$1$ (Theorem~\ref{thm:fixed} and
  Proposition~\ref{prop:det}). For a nondegenerate quaternion subalgebra
  $Q=\spn\{1,x,y,xy\}$ this places the defect in $\SO(Q^\perp)$.
\item A census over the split octonions modulo small primes supports an
  exact rank law (Conjecture~\ref{conj:rank}). Its first case is proved
  (Proposition~\ref{prop:flat}).
\item A Lean~4 oracle recomputes the five edges of the associator pentagon
  from an explicit structure tensor. A mutation analysis
  (Table~\ref{tab:mutants}) shows which checks detect which implementation
  errors, and it identifies two error classes that closure of the pentagon
  cannot detect (Observation~\ref{obs:weak}).
\end{enumerate}

\subsection*{Scope and claims}
We make no novelty claim for the algebra. Proposition~\ref{prop:basic},
Theorem~\ref{thm:fixed} and Proposition~\ref{prop:det} follow quickly from
the composition identities and Artin's theorem, and may well appear in the
literature in some form. A dated prior-art search has not been completed (see
Section~\ref{sec:related}). The computations are finite and seeded. They
support Conjecture~\ref{conj:rank} but do not prove it, and the census is
evidence, not a theorem. The Lean component is an executable checker, not a
kernel-checked proof of soundness.

\section{Preliminaries}\label{sec:prelim}

Throughout, $k$ is a field with $\operatorname{char}k\neq2$. A
\emph{composition algebra} over $k$ is a unital $k$-algebra $\C$ with a
nondegenerate quadratic form $N$ satisfying $N(xy)=N(x)N(y)$. We write
$B(x,y)=N(x+y)-N(x)-N(y)$ for the polar form, $T(x)=B(x,1)$ for the trace,
and $\conj{x}=T(x)1-x$ for the canonical involution. An \emph{octonion
algebra} is an eight-dimensional composition algebra. We use the following
classical facts~\cite{schafer1966,springer2000}.

\begin{enumerate}[label=(F\arabic*)]
\item\label{F:comp} $\C$ is alternative, $x\conj{x}=\conj{x}x=N(x)1$, and
  \[
    x(\conj{x}w)=\conj{x}(xw)=N(x)w\qquad(x,w\in\C).
  \]
  In particular $L_x$ is invertible exactly when $N(x)\neq0$, with
  $L_x^{-1}=N(x)^{-1}L_{\conj{x}}$.
\item\label{F:artin} \emph{Artin's theorem.} In an alternative algebra, the
  subalgebra generated by any two elements is associative.
\end{enumerate}

The associator is $\assoc{x}{y}{z}=x(yz)-(xy)z$. It is trilinear and vanishes
whenever one argument is~$1$.

\subsection*{The finite model}
Let $\ell$ be an odd prime. Iterating the Cayley--Dickson doubling
\begin{equation}\label{eq:cd}
  (a,b)(c,d)=(ac+\gamma\,\conj{d}b,\ da+b\conj{c}),\qquad
  \conj{(a,b)}=(\conj{a},-b),
\end{equation}
three times over $\F_\ell$ with $\gamma=-1$ at each stage gives an octonion
algebra $\C_\ell$ on $\F_\ell^8$ with $N(x)=\sum_i x_i^2$. The diagonal form
in eight variables is isotropic over $\F_\ell$ by the Chevalley--Warning
theorem~\cite{chevalley1935,warning1935}. Hence $\C_\ell$ is the split
octonion algebra, which is unique up to isomorphism (Zorn's vector-matrix
algebra~\cite{zorn1933}; see also~\cite[Ch.~1]{springer2000}). Our
computations use this model; the theory of Section~\ref{sec:defect} applies to
every octonion algebra over every field of characteristic $\neq2$.

\section{The square-root-free defect}\label{sec:defect}

\begin{definition}
For $x,y\in\C$ with $N(x)N(y)\neq0$, set
\[
  \Delta(x,y)=L_xL_yL_{xy}^{-1}
  =\bigl(N(x)N(y)\bigr)^{-1}L_xL_yL_{\conj{xy}}.
\]
The second expression follows from~\ref{F:comp}, since
$N(xy)=N(x)N(y)\neq0$. Set $A=A(x,y)$ for the unital subalgebra generated by
$x$ and $y$, and $Q=Q(x,y)=\spn\{1,x,y,xy\}\subseteq A$.
\end{definition}

\begin{remark}\label{rem:real}
For $k=\R$ and $\C$ the division octonions,
$U_xU_yU_{xy}^{-1}=\frac{\sqrt{N(xy)}}{\sqrt{N(x)}\sqrt{N(y)}}\,\Delta(x,y)
=\Delta(x,y)$. Thus $\Delta$ extends~\eqref{eq:real-defect}, and it uses only
field operations.
\end{remark}

\begin{proposition}\label{prop:basic}
Let $N(x)N(y)\neq0$ and $\Delta=\Delta(x,y)$. Then
\begin{enumerate}[label=\textup{(\alph*)}]
\item $\Delta\in\Ort(\C,N)$;
\item $\Delta(xy)=xy$;
\item $\Delta=I$ if and only if $\assoc{x}{y}{z}=0$ for every $z\in\C$.
\end{enumerate}
\end{proposition}

\begin{proof}
(a) By composition, $N(L_av)=N(a)N(v)$ for all $a,v$. Applying this to
$v=L_{xy}^{-1}u$ gives $N(u)=N(xy)\,N(L_{xy}^{-1}u)$. Hence
$N(\Delta u)=N(x)N(y)N(xy)^{-1}N(u)=N(u)$, and polarization gives
invariance of~$B$.
(b) $L_{xy}^{-1}(xy)=1$, so $\Delta(xy)=x(y\cdot1)=xy$.
(c) $\Delta=I$ if and only if $L_xL_y=L_{xy}$, and evaluating at every $z$ gives
$x(yz)=(xy)z$.
\end{proof}

\begin{theorem}\label{thm:fixed}
Let $N(x)N(y)\neq0$ and $\Delta=\Delta(x,y)$. Then:
\begin{enumerate}[label=\textup{(\alph*)}]
\item $\Delta w=w$ for every $w\in A$, and in particular for every $w\in Q$;
\item $\Delta(W^\perp)=W^\perp$ for every subspace $W\subseteq A$;
\item if $B|_W$ is nondegenerate for a subspace $W\subseteq A$, then
  $\C=W\oplus W^\perp$ and $\Delta=\mathrm{id}_W\oplus\Delta|_{W^\perp}$ with
  $\Delta|_{W^\perp}\in\Ort(W^\perp,N)$;
\item $\rank(\Delta-I)\le 8-\dim A\le 8-\dim Q$.
\end{enumerate}
\end{theorem}

\begin{proof}
(a) Let $A_0$ be the subalgebra generated by $x$ and $y$ without the unit. By
\ref{F:comp} and \ref{F:artin}, $A_0$ is associative. Every element of $A$ has
the form $\alpha1+a_0$ with $a_0\in A_0$. Since the associator is trilinear and
vanishes when an argument is $1$, it vanishes on $A$, so $A$ is associative.
Let $w\in A$ and put $z=N(xy)^{-1}\conj{xy}\,w$. Since
$\conj{xy}=T(xy)1-xy\in A$, we have $z\in A$, and
$z=L_{xy}^{-1}w$ by~\ref{F:comp}. Associativity of $A$ gives
\[
  \Delta w=x(yz)=(xy)z=L_{xy}L_{xy}^{-1}w=w.
\]
(b) Let $u\in W^\perp$ and $w\in W$. By (a) and
Proposition~\ref{prop:basic}(a),
$B(\Delta u,w)=B(\Delta u,\Delta w)=B(u,w)=0$, so
$\Delta(W^\perp)\subseteq W^\perp$. Equality holds because $\Delta$ is injective
and $W^\perp$ is finite-dimensional.
(c) Nondegeneracy of $B|_W$ gives $W\cap W^\perp=0$ and
$\dim W^\perp=8-\dim W$. The rest follows from (a), (b) and
Proposition~\ref{prop:basic}(a).
(d) By (a), $A\subseteq\ker(\Delta-I)$.
\end{proof}

\begin{proposition}\label{prop:det}
For every $x\in\C$, $\det L_x=N(x)^4$. Consequently $\det\Delta(x,y)=1$
whenever $N(x)N(y)\neq0$, so $\Delta(x,y)\in\SO(\C,N)$.
\end{proposition}

\begin{proof}
Fix a basis $e_0,\dots,e_7$ of $\C$, let $R=k[t_0,\dots,t_7]$, and let
$\xi=\sum_it_ie_i\in\C_R=\C\otimes_kR$. The entries of $L_\xi$ and
$L_{\conj\xi}$ are linear forms in the $t_i$. Expanding,
\[
  \conj\xi(\xi w)
  =\sum_it_i^2\,\conj{e_i}(e_iw)
   +\sum_{i<j}t_it_j\bigl(\conj{e_i}(e_jw)+\conj{e_j}(e_iw)\bigr).
\]
By~\ref{F:comp} and its polarization in $\C$, the right-hand side is
$\bigl(\sum_it_i^2N(e_i)+\sum_{i<j}t_it_jB(e_i,e_j)\bigr)w=N(\xi)w$. Hence
$L_{\conj\xi}L_\xi=N(\xi)I$ over $R$, and
\[
  \det L_{\conj\xi}\cdot\det L_\xi=N(\xi)^8\quad\text{in }R.
\]
A nondegenerate quadratic form in at least three variables is irreducible: a
product of two linear forms has rank at most~$2$. Since $R$ is a unique
factorization domain, $\det L_\xi=c\,N(\xi)^a$ with $c\in k^\times$. As
$\det L_\xi$ is homogeneous of degree $8$, we get $a=4$, and specializing
$\xi\mapsto1$ gives $c=\det I=1$. Specializing $t$ gives $\det L_x=N(x)^4$ for
every $x\in\C$. Finally,
\[
  \det\Delta(x,y)=\bigl(N(x)N(y)\bigr)^{-8}\,N(x)^4N(y)^4N(\conj{xy})^4=1,
\]
since $N(\conj{xy})=N(xy)=N(x)N(y)$.
\end{proof}

\begin{corollary}\label{cor:so4}
If $\dim Q=4$ and $B|_Q$ is nondegenerate, then
$\Delta(x,y)=\mathrm{id}_Q\oplus R$ with $R\in\SO(Q^\perp,N)$ and $\rank(\Delta-I)\le4$.
\end{corollary}

\begin{proof}
By Theorem~\ref{thm:fixed}(c) with $W=Q$, $R=\Delta|_{Q^\perp}$ is an isometry.
By Proposition~\ref{prop:det}, $\det R=\det\Delta=1$.
\end{proof}

\begin{remark}
Proposition~\ref{prop:basic} places the defect in the stabilizer of $xy$, a
copy of $\Ort(7)$. Corollary~\ref{cor:so4} shrinks this to the pointwise
stabilizer of a four-dimensional subspace, a copy of $\SO(4)$. Over
$\F_\ell$ the latter is a group of order $O(\ell^6)$, compared with
$O(\ell^{21})$ for the former.
\end{remark}

\begin{conjecture}[rank law]\label{conj:rank}
Let $\C$ be the split octonion algebra over $\F_q$, $q$ odd, and let
$N(x)N(y)\neq0$. Then $\rank(\Delta(x,y)-I)=2$ if $\dim Q=3$ and $\rank(\Delta(x,y)-I)=4$ if $\dim Q=4$. Together with Proposition~\ref{prop:flat}, this says
\[
  \rank\bigl(\Delta(x,y)-I\bigr)=
  \begin{cases}0,&\dim Q\le2,\\ 2,&\dim Q=3,\\ 4,&\dim Q=4.\end{cases}
\]
\end{conjecture}

The first case is a theorem (Proposition~\ref{prop:flat}), and
Theorem~\ref{thm:fixed}(d) gives the upper bound in the third. The evidence for
the remaining equalities is the census in Section~\ref{sec:census}.

\begin{proposition}\label{prop:flat}
If $\dim Q\le2$, then $\Delta(x,y)=I$.
\end{proposition}

\begin{proof}
If $\dim Q\le2$, then $1,x,y$ are linearly dependent. If $x\in k1$, then $L_x$ is
scalar and $\assoc xyz=0$ for every $z$. Otherwise $y=\alpha1+\beta x$ for some
$\alpha,\beta\in k$, and
$\assoc xyz=\beta\assoc xxz=0$ by alternativity~\ref{F:comp}. In either case
$\Delta=I$ by Proposition~\ref{prop:basic}(c).
\end{proof}

\section{The associator pentagon}\label{sec:pentagon}

Let $a,b,c,d$ be elements of an arbitrary (not necessarily associative)
algebra over a commutative ring. The five binary parenthesizations of $abcd$,
\[
  v_0=((ab)c)d,\quad v_1=(a(bc))d,\quad v_2=a((bc)d),\quad
  v_3=a(b(cd)),\quad v_4=(ab)(cd),
\]
form the vertices of a pentagon, the two-dimensional associahedron
\cite{stasheff1963,maclane1963}. The edge differences $e_i=v_{i+1}-v_i$
(indices modulo~$5$) are associators:
\begin{equation}\label{eq:edges}
  e_0=\assoc abc\,d,\quad e_1=\assoc{a}{bc}{d},\quad e_2=a\,\assoc bcd,\quad
  e_3=-\assoc{a}{b}{cd},\quad e_4=-\assoc{ab}{c}{d}.
\end{equation}
Since $\sum_ie_i$ telescopes,
\begin{equation}\label{eq:teich}
  \assoc abc\,d+\assoc{a}{bc}{d}+a\,\assoc bcd-\assoc{a}{b}{cd}-\assoc{ab}{c}{d}=0,
\end{equation}
which is the Teichm\"uller identity~\cite{zhevlakov1982}. In a
nonassociative algebra the individual edges may be nonzero, so the identity
records non-associativity as edge data on a closed boundary.

\begin{observation}\label{obs:weak}
As a test of an implementation of~\eqref{eq:edges}, the closure
\eqref{eq:teich} is weak in two specific ways.
\begin{enumerate}[label=\textup{(\alph*)}]
\item It holds in every bilinear algebra. An implementation that computes all
  products in the wrong algebra therefore still closes, for example the
  opposite algebra obtained by transposing the structure tensor.
\item Replacing $\assoc xyz$ by $-\assoc xyz$ negates every edge and preserves
  the closure.
\end{enumerate}
Detecting (a) requires comparison with an independent computation in the
intended algebra. Detecting (b) requires checking each edge against its
vertex difference $v_{i+1}-v_i$, which is computed without associators.
\end{observation}

\section{Computational certification}\label{sec:census}

\subsection{Defect census}
For each $\ell\in\{3,5,7,11\}$ we draw $400$ pairs $(x,y)$ of nonzero-norm
elements of $\C_\ell$ uniformly with a fixed seed. We add degenerate pairs
$y=\alpha1+\beta x$ ($\beta\neq0$, $N(y)\neq0$) for a quarter of the $x$.
Each pair is checked exactly in $\F_\ell$ for:
\begin{itemize}
\item the inverse formula;
\item Proposition~\ref{prop:basic}(a)--(c);
\item $\det\Delta=1$;
\item $\Delta|_Q=\mathrm{id}$;
\item the rank law.
\end{itemize}
Table~\ref{tab:census} reports the strata. No check failed. On the
$\dim Q=4$ stratum, $\tr\Delta$ takes every value in $\F_\ell$, so the trace is
not determined by $\dim Q$. Whether the class of $\Delta|_{Q^\perp}$ carries
arithmetic information beyond classical conjugacy data is the open question of
Section~\ref{sec:open}. The stratum $\dim Q=3$ is rare (six pairs in total)
and absent for $\ell\ge7$, so the middle case of Conjecture~\ref{conj:rank}
rests on little evidence.

\begin{table}[t]
\centering
\small
% Generated by paper/associator-defects/make_tables.py from data/n2/defect-census-v1.json; do not edit.

\begin{tabular}{rrcccrr}
\toprule
 & & \multicolumn{3}{c}{$(\dim Q,\ \operatorname{rank}(\Delta-I))$} & distinct $\operatorname{tr}\Delta$ & \\
\cmidrule(lr){3-5}
$\ell$ & pairs & $(2,0)$ & $(3,2)$ & $(4,4)$ & at $\dim Q=4$ & violations \\
\midrule
3 & 481 & 82 & 4 & 395 & 3 & 0 \\
5 & 480 & 80 & 2 & 398 & 5 & 0 \\
7 & 484 & 84 & 0 & 400 & 7 & 0 \\
11 & 487 & 87 & 0 & 400 & 11 & 0 \\
\bottomrule
\end{tabular}

\caption{Defect census over $\C_\ell$: number of pairs in each stratum
$(\dim Q,\rank(\Delta-I))$, number of distinct traces on the $\dim Q=4$
stratum, and the total number of failed checks. Generated from
\texttt{data/n2/defect-census-v1.json}.}
\label{tab:census}
\end{table}

\subsection{A Lean oracle for the pentagon}
The oracle, written in Lean~4~\cite{moura2021}, takes a modulus
$p>1$, an explicit structure tensor $t_{ijk}$ (the coefficient of $e_k$ in
$e_ie_j$) and four vectors. It recomputes every product from the tensor
alone. Its self-check runs on two models:
\begin{itemize}
\item an $\F_3$ octonion tensor with a basis witness;
\item an $\F_5$ Cayley--Dickson tensor generated, together with three input
  quadruples and their edge values, by an independent Python implementation
  of~\eqref{eq:cd}.
\end{itemize}
Two of the $\F_5$ quadruples have all five edges nonzero.

For each $\F_5$ case the self-check requires the following:
\begin{itemize}
\item \emph{closure}: \eqref{eq:teich} holds;
\item \emph{replay}: the five edges equal the independently computed values;
\item \emph{vertex}: each edge equals its vertex difference;
\item \emph{sign flip}: negating edge $i$ breaks the closure exactly when
  $e_i\neq0$ (true for odd $p$, since the closure then changes by $2e_i$);
\item in cases with all edges nonzero, moving the outer factor of
  $\assoc abc d$ or $a\assoc bcd$ to the other side breaks the closure.
\end{itemize}

Table~\ref{tab:mutants} lists six single-fragment mutants of the oracle source
and the check families that detect each one. Every mutant is detected. The
two classes of Observation~\ref{obs:weak} behave as predicted:
\begin{itemize}
\item the globally negated associator preserves closure, and is detected only
  by replay, the vertex check, and the pinned $\F_3$ witness;
\item the transposed tensor (the opposite algebra) passes closure and the
  vertex check, and is detected \emph{only} by replay against the independent
  implementation.
\end{itemize}
The basis quadruples $(e_0,e_1,e_2,e_4)$ of both models, whose first factor is
the identity, detect none of the four single-edge mutants. Inputs with every edge
nonzero are necessary.

\begin{table}[t]
\centering
\small
% Generated by paper/associator-defects/make_tables.py from data/n2/oracle-mutants-v1.json; do not edit.

\begin{tabular}{lcccccr}
\toprule
Mutant of \texttt{Oracle.lean} & Closure & Replay & Vertex & Sign flip & $\mathbb F_3$ & failing checks \\
\midrule
$e_3 = [a,b,cd]$ (sign dropped) & $\bullet$ & $\bullet$ & $\bullet$ & $\bullet$ &  & 8 \\
$e_0 = d\,[a,b,c]$ (factor side) & $\bullet$ & $\bullet$ & $\bullet$ &  &  & 6 \\
$e_1 = [bc,a,d]$ (argument order) & $\bullet$ & $\bullet$ & $\bullet$ & $\bullet$ &  & 8 \\
$[x,y,z] = (xy)z - x(yz)$ (global sign) &  & $\bullet$ & $\bullet$ &  & $\bullet$ & 8 \\
$e_4 = -[ba,c,d]$ (inner order) & $\bullet$ & $\bullet$ & $\bullet$ &  &  & 6 \\
$t_{ijk} \mapsto t_{jik}$ (opposite algebra) &  & $\bullet$ &  &  &  & 2 \\
\bottomrule
\end{tabular}

\caption{Mutation analysis of the Lean oracle: check families that report
failure for each mutant, and the number of failing named checks. Generated by
\texttt{scripts/lean\_mutation\_suite.py}
(\texttt{data/n2/oracle-mutants-v1.json}).}
\label{tab:mutants}
\end{table}

\subsection{Reproducibility}
All data files are generated deterministically and checked byte for byte in
continuous integration. One portability defect is worth recording. Since
version~3.12, Python's built-in \texttt{sum} uses compensated summation over
floats, so descriptive statistics computed under 3.11 and 3.12 differ in the
last bits. A statistic that is mathematically zero can come out as $0$ under
one version and as $\approx10^{-31}$ under the other. Committed
floating-point summaries are therefore rounded to $10$ decimal places and
$12$ significant digits before hashing. The exact integer data (orbit counts,
strata, edges) are unaffected.

Claims are recorded in a content-addressed ledger of proof records, which
distinguishes verified finite computations, repository theorems, imported
theorems, bounded experiments and pending statements. A claim-governance
check keeps the ledger consistent with its rendered surfaces and requires a
test for every claim marked proved. Appendix~\ref{app:repro} lists the
artifacts.

\section{Related work}\label{sec:related}

Composition algebras and their classification go back to
Hurwitz~\cite{hurwitz1898}. Standard references are Schafer~\cite{schafer1966}
and Springer--Veldkamp~\cite{springer2000}, and the surveys of
Baez~\cite{baez2002} and Conway--Smith~\cite{conway2003} cover normalized left
multiplications, triality and $\mathrm{Spin}(8)$. Artin's theorem and the
Teichm\"uller identity are classical~\cite{schafer1966,zhevlakov1982}. The
pentagon of parenthesizations is the associahedron of
Stasheff~\cite{stasheff1963} and the coherence pentagon of
Mac~Lane~\cite{maclane1963}. Here it appears in its simplest, linear form.
Mutation analysis as a measure of test adequacy goes back to DeMillo, Lipton
and Sayward~\cite{demillo1978}.

We expect that statements equivalent to Theorem~\ref{thm:fixed} and
Proposition~\ref{prop:det} are known to specialists in some form, possibly
phrased via triality or the Moufang identities. A systematic, dated search of
the literature on triality, octonionic loops, and defect or holonomy
constructions has not yet been carried out. Until it is, the paper claims
only the exact finite formulation, the explicit proofs, and the certification
methodology.

\section{Limitations and open problems}\label{sec:open}

\begin{enumerate}[label=(\arabic*)]
\item \emph{Rank law.} Prove or refute Conjecture~\ref{conj:rank}. The case
  $\dim Q=3$ needs targeted rather than uniform sampling.
\item \emph{Defect-stratified statistics.} For navigation words in generators
  drawn from a prime-norm shell modulo $\ell$, do collision moments
  conditioned on the conjugacy class of $\Delta|_{Q^\perp}$ in $\SO(Q^\perp)$
  differ from the unconditioned moments, beyond what conjugacy data in
  $\Ort(8,\F_\ell)$ predicts? A negative answer would show that the defect is
  classical bookkeeping only.
\item \emph{Formal soundness.} The oracle decides finite instances by
  evaluation. A kernel-checked theorem connecting its Boolean checks to the
  composition-algebra structures of Lean's mathematical library is not yet
  available.
\item \emph{Scope of the census.} Sample sizes are small and seeded. The
  census establishes no asymptotic or distributional statement.
\end{enumerate}

\appendix
\section{Reproduction}\label{app:repro}

The repository \texttt{larsbx/novelty-lab} contains every artifact. With
Python~$\ge3.11$ and the Lean toolchain \texttt{leanprover/lean4:v4.24.0}:
\begin{center}
\small
\begin{tabular}{ll}
\toprule
Artifact & Command \\
\midrule
Census (Table~\ref{tab:census}) & \texttt{python experiments/n2/defect\_census.py --check} \\
Theorem replay & \texttt{python -m unittest discover -s tests -p 'test\_n2\_defect*.py'} \\
Lean oracle & \texttt{lake build \&\& lake exe oracle-selfcheck} \\
Mutants (Table~\ref{tab:mutants}) & \texttt{python scripts/lean\_mutation\_suite.py} \\
Paper tables & \texttt{python paper/associator-defects/make\_tables.py --check} \\
\bottomrule
\end{tabular}
\end{center}
Census seed \texttt{20261001}. SHA-256 digests:
\begin{itemize}
\item census data: \texttt{411511f2\ldots8eb83e};
\item mutant data: \texttt{bac544cb\ldots9af2fe7};
\item generated Lean fixture: \texttt{3d07bac0\ldots4377898};
\item oracle source (unmutated): \texttt{a19ee3dd\ldots ed86e90e}.
\end{itemize}

\section*{Acknowledgements}
Code, computations and a first draft of this manuscript were produced with the
assistance of an AI coding agent (Claude Code). The authors take
responsibility for all statements, proofs and data, and will verify every
proof and citation before submission.

\bibliographystyle{amsplain}
\bibliography{references}

\end{document}
